\setchapterpreamble[u]{\margintoc}
\chapter{Indefinite object drop} % Object-dropping verbs
\labch{objectdrop}

As mentioned in \refch{intro}, this thesis is about indefinite implicit objects. What is "indefinite" about them? In what sense can they be considered "implicit"? And ultimately, what is objecthood itself? This Chapter will answer these questions in reverse order, from the most general to the most specific one. In \refsec{theory_transitivity} I will make reference to the transitivity continuum, as defined by \textcite{HopperThompson1980} and further explored by later literature. In \refsec{theory_def_vs_indef} a crucial distinction will be made between definite and indefinite object drop, following \textcite{Fillmore1986} and subsequent works. The nature of \textit{indefinite} object drop will be finally described in \refsec{theory_defindefinite} and \refsec{theory_entries}, and a working definition (for the purposes of this thesis) will be provided in \refsec{theory_workingdef}.

\section{Transitivity as a prototype} \labsec{theory_transitivity}

% \subsection{Transitivity in Hopper \& Thompson (1980)} \labsec{theory_ht1980}
School kids everywhere are used to call "transitive" the verbs which take an overt direct object, with reference to the syntax. In a traditional semantic definition, a clause is deemed "transitive" if it describes an event where the action performed by an Agent "passes over"\sidenote{Hence the name of the concept of transitivity, from Latin \textit{transire} 'to go over'.} to a Patient, which usually undergoes some kind of transformation.\\
Going beyond these naive definitions, but still capturing their spirit, \textcite{HopperThompson1980} first proposed an account where transitivity is interpreted as a scalar concept whose strength depends on several parameters, or, to use more modern terminology, as a prototype category \parencite{Naess2007}. In particular, they identified ten parameters \parencite[252]{HopperThompson1980}, reported almost \textit{verbatim} in \reftab{ht1980_parameters}.

\begin{table}[htb] % the "htb" makes table env unfloaty
\caption{\textcite[252]{HopperThompson1980} defined transitivity as a prototype concept determined by ten parameters.}
\labtab{ht1980_parameters}
\begin{tabular}{rl|ll}
 & & \textbf{high transitivity} & \textbf{low transitivity} \\
 \hline
A. & \textbf{Participants} & 2+ (Agent and Object)  & 1 participant  \\
B. & \textbf{Kinesis} & action  & non-action  \\
C. & \textbf{Aspect} & telic  & atelic  \\
D. & \textbf{Punctuality} & punctual  & non-punctual  \\
E. & \textbf{Volitionality} & volitional  & non-volitional  \\
F. & \textbf{Affirmation} & affirmative  & negative  \\
G. & \textbf{Mode} & realis  & irrealis  \\
H. & \textbf{Agency} & A high in potency  & A low in potency   \\
I. & \textbf{Affectedness of O} & O totally affected  & O not affected  \\
J. & \textbf{Individuation of O} & O highly individuated  & O non-individuated  
\end{tabular}
\end{table}

These parameters are potentially active in all languages, but languages may differ from one another with respect to the actual subset of parameters they select as necessary criteria for transitivity. This depends on the "recursivity" \parencite[29]{Naess2007} of prototypical concepts, which assign membership in a category (in this case, transitive clauses) on the basis of attributes which are prototype concepts themselves \parencite[61]{taylor1995linguistic}.\\
Parameters A, B, E, F, and G from \reftab{ht1980_parameters} are self-explanatory. Parameter C (telicity) will be discussed in more detail in \refsec{telicity}. Parameter D (punctuality) refers to the phase between inception and completion of an action, which is non-existent with verbs like \textit{to kick} and noticeable with verbs like \textit{to carry}. Parameter H (agency) separates animate and inanimate subjects. Parameter I (affectedness of the object) determines that sentences like \textit{I drank up the milk} are more transitive than sentences like \textit{I drank some of the milk}, since the milk is only partially affected by the drinking in the second sentence. Finally, parameter J (individuation of the object) refers to the distinctness of the object both from the Agent and from the background, as summarized in \reftab{ht1980_individuation} \parencite[253]{HopperThompson1980}.

\begin{table}[htb] % the "htb" makes table env unfloaty
\caption{\textcite[252]{HopperThompson1980} defined transitivity as a prototype concept determined by ten parameters.}
\labtab{ht1980_individuation}
\begin{tabular}{c|c}
 \textbf{individuated} & \textbf{non-individuated} \\
 \hline
proper & common \\
human, animate & inanimate \\
concrete & abstract \\
singular & plural \\
count & mass \\
referential, definite & non-referential
\end{tabular}
\end{table}

The individuation parameter is the most controversial among the ten proposed ones, as observed by \textcite[128]{comrie1989language} and later on by \textcite[18]{Naess2007}. According to what has come to be known as "Comrie's generalization", in prototypical transitive clauses both animacy and definiteness are high in the Agent and low in the Patient, \textit{contra} \textcite{HopperThompson1980}. The weak argumentation \textcite{HopperThompson1980} provide in favor of the individuation parameter is that speakers would be more likely to focus on the Patient in \textit{I bumped into Charles} than in \textit{I bumped into the table}, since bumping is more likely to affect human beings than tables. \textcite{comrie1989language} makes a much more compelling point with reference to cross-linguistic typology, basing the generalization on the animacy hierarchy and on referential case-marking (which I will not discuss here, since it would lead me too far from the intended argumentation). Later literature \parencite{Naess2007, kemmer1993middle, Kardos2010, Naess2009} reinforced this point by assuming that prototypical transitive events are described by verbs whose subject and object are maximally distinct from a semantic point of view.\\
To sum up, the terse summary by \textcite[15]{Naess2007} clearly shows the relation between the naive definitions of transitivity and the ten-parameter account by \textcite{HopperThompson1980}. A prototypical clause is understood to describe an event such that:
\begin{itemize}
    \item a volitional Agent (E, H)
    \item performs an action (B)
    \item with a tangible, lasting effect on a Patient (A, I, J),
    \item and it is presented as real and concluded (C, D, F, G).
\end{itemize}

\textcite[78]{Lorenzetti2008} provides a tighter cluster of parameters, arguing that only a subset of the ones proposed by \textcite{HopperThompson1980} are truly relevant. In particular, the author ditches the criteria relative to the transitive event being real and concluded (C, D, F, G), and only keeps agentivity, affectedness, and individuation of the object among the other groups of criteria.


\section{Definite \textit{vs} indefinite drop} \labsec{theory_def_vs_indef}

In \refsec{theory_transitivity} I introduced transitivity as a prototype concept depending on a cluster of parameters and, specifically, involving an Agent acting on a Patient. But what about the utterances where events of this kind are expressed without an overt syntactic object? In this Section I will provide an account of the literature on the matter.

\subsection{Either definite or indefinite: discrete accounts} \labsec{theory_discrete}

\paragraph{Introduction}
Verbs behaving intransitively are, using \textcite[191]{rutherford1998workbook}'s words, "a mixed bag". Consider, for instance, the examples in \ref{intro}. The sentence in \ref{intro1} features a typical intransitive verb, describing an event where the subject is \textit{not} performing an action with effects on some Patient. The sentence in \ref{intro2} also describes an event where the subject is not volitionally acting on a Patient, but it is nevertheless clear that there has to exist something that John knows (unlike in \ref{intro1}, where there cannot be something that John sleeps). The sentence in \ref{intro3} has an Agent acting volitionally on a Patient, which is however not instantiated syntactically as an overt direct object.

\ex. \label{intro} \a. \label{intro1} John slept.
\b. \label{intro2} John knew.
\c. \label{intro3} John ate.

There is a clear similarity between \ref{intro2} and \ref{intro3}, as opposed to \ref{intro1}. They both require a Patient/Theme semantically \parencite[510]{Somers1984}, and they both surface as object-less syntactically. Quoting \textcite[48]{Yasutake1987}, "they are different from pure intransitives in that the action will not be complete without some lexically implied (but unspecified) object". Oddly, some literature (\textcite{BourmayanRecanati2013, Liu2008} a.o.) does not interpret such verbs as transitive-become-intransitive verbs via the omission of the direct object, but as intransitive-made-transitive verbs. Such an interpretation is totally counter-intuitive and it goes against the generally-accepted tenet that a core feature of so-called "intransitive verbs" is that they have no object slot available in the syntax.\\
There is, however, a crucial difference between these \ref{intro2} and \ref{intro3}. Native speakers of English understand that they have to be provided some context in order to understand \ref{intro2} \textemdash what is it exactly that John knew? On the contrary, \ref{intro3} can be interpreted to mean that John had a meal at a certain moment in time, without resorting to any additional context. This distinction was captured and defined by \textcite{Fillmore1986} (building upon \textcite{fillmore1969types, Allerton1975}), the established seminal work on the distinction between so-called "definite" and "indefinite" object drop (here represented by \ref{intro2} and \ref{intro3}, respectively).

\paragraph{Fillmore's account}
\textcite[96]{Fillmore1986} distinguishes between Indefinite Null Complements (hence, INC) and Definite Null Complements (hence, DNC) by testing "whether it would sound odd for a speaker to admit ignorance of the identity of the reference of the missing phrase". So, making reference to \ref{intro} again, there would be no issue with saying "John ate. I wonder what he ate.", but it would be quite odd to say "John knew. I wonder what he knew.". Thus, the missing object in \ref{intro2} is a DNC, while the missing object in \ref{intro3} is an INC. Fillmore than splits INCs into two sub-groups based on whether the omitted object is "of considerable generality" or "requiring the specification of various degrees of semantic specialization". The examples in \ref{fillmore} \parencite[96-97]{Fillmore1986} show increasing degrees of, using his words, "semantic specialization". In \ref{fillmore1}, the subject cannot perform the very act of eating or drinking, regardless of the actual ingested item. In \ref{fillmore2} something specific was eaten, but its specific nature is irrelevant inasmuch the speaker is referring to eating as the act of having a meal. In \ref{fillmore3} the omitted object is referring not just to any drinkable liquid, but to alcohol specifically. Finally, in \ref{fillmore4} something very specific was baked by the subject, but this information is backgrounded to focus on the activity itself (I will come back to this in \refsec{theory_incorporation}).

\ex. \label{fillmore} \a. \label{fillmore1} When my tongue was paralyzed I couldn't eat or drink.
\b. \label{fillmore2} We've already eaten.
\c. \label{fillmore3} I've tried to stop drinking.
\d. \label{fillmore4} I spent the afternoon baking.

However, as it will be shown throughout this Chapter, this secondary division of INCs into subgroups does not have to be a binary theoretical distinction. On the contrary, it follows from several finer-grained considerations on the nature of INCs and the factors allowing them. Moreover, this binary division actually opens the door to discussions about a DNC-INC continuum (more details on this in \refsec{theory_continuous}). Specifically, how is "semantic specialization" different from the "knownness" of the object in DNC constructions \parencite[525]{Eu2018}? An answer can be found in \textcite[218]{Allerton1975}, where the case is made that semantically specialized INCs (just like any INC) refer to a category of individuals, while DNCs refer to specific instances of a given category.\\
Going back to the main distinction between INCs and DNCs, finally, \textcite{Fillmore1986} formally defines the former as objects whose "referent's identity is unknown or a matter of indifference" and the latter as objects whose referent "must be retrieved from something \textit{given} in the context". This context "has either to be given linguistically, in the preceding context, or extralinguistically, in the situational context" \parencite[13]{StarkMeier2017}.

\paragraph{Other accounts}
Several other researchers made use of the distinction between definite and indefinite implicit objects brought to the fore by \textcite{Fillmore1986}, providing slightly different definitions which capture different aspects of the phenomenon. \textcite{Allerton1975}, then further developed by Fillmore, distinguishes between "contextual omission" and "lexical omission", which respectively refer to Fillmore's DNCs and INCs. \textcite{CumminsRoberge2004} distinguish between "internally-licensed null objects" (INCs) and "referential null objects" (DNCs). According to \textcite[30]{PethoKardos2006}, INCs "receive an \textit{existentially quantified} interpretation", while DNCs are "interpreted anaphorically and must therefore have an appropriate antecedent in context to make sense" (\textcite{Fillmore1986} himself referred to this in the title itself, "Pragmatically Controlled Zero Anaphora"). \textcite[13]{Medina2007}, the foundational work upon which I am basing my own model of the indefinite object construction, sees DNCs as "implicit objects whose particular meaning can be recovered from the preceding discourse or disambiguating physical context" and INCs as "implicit objects whose meaning is recoverable only from the verb in the sentence". Here the focus is all on recoverability, and the author goes on to show that semantic recoverability can be a reliable predictor of object drop in INC sentences. \textcite[293]{Liu2008}, following \textcite{Garcia-VelascoMunoz2002}, takes the shift away from lexical semantics and onto aspectual territory. In particular, the point is made that INCs involve a change of focus "from the object in the transitive use to the activity (the verb) itself in the intransitive use" (an idea that I will discuss in full detail in \refsec{theory_incorporation}), while DNCs do not determine such a shift.\\
The accounts provided so far are not in conflict with Fillmore's formulation of the problem at hand, nor are they in conflict with the view I am adopting in this thesis in order to provide a probabilistic model of the implicit object construction. Other accounts, on the other hand, are more challenging and deserve a more thorough clarification. Let us consider the most relevant ones for my argumentation.\\
\textcite[55]{TonelliDelmonte2011} argue that, while INCs are "constructionally licensed, in that they apply to any predicate in a particular grammatical construction", DNCs are "lexically specific, in that they apply only to some predicates". Later in this Section (on \refpage{recipes}) and in \refsec{theory_continuous} I will bring evidence in support of the opposite point of view, which is in favor of seeing DNCs as (extra- and intra-linguistically) contextually, not lexically, determined. Moreover, \textcite{TonelliDelmonte2011}'s account is in direct conflict with \textcite[95]{Fillmore1986}, who argues that INCs are "limited to particular lexically defined environments" (such as the object slot of \textit{to eat} and \textit{to read}). In this regard, I side with \textcite{TonelliDelmonte2011}. My probabilistic model of INCs (the results thereof are discussed in \refch{results} and modeled in \refch{model}) will provide strong evidence in support of the idea that any transitive verb can participate in INC constructions, provided the right aspectual, semantic, and pragmatic features. Indeed, as noted by \textcite[216]{HuddlestonEtAl2002}, transitivity is better thought of as a property of verb \textit{use}, rather than a feature of verbs themselves.\\
In a pragmatic (in particular, not lexical) perspective, \textcite[44]{AnderBois} and \textcite[53-54]{Melchin2019} both stay true to Fillmore's original interpretation of DNCs as "pragmatic anaphoras", arguing that DNCs corefer with other referents in the discourse. Moreover, they maintain that INCs lack the possibility of having coreferential interpretations. Fillmore himself \parencite[97]{Fillmore1986} made this point with example \ref{coref}, where \ref{coref2} cannot be considered a proper answer to the question in \ref{coref1}. This is taken to mean that there is no co-reference between the sandwich and the implicit object of \textit{to eat} in \ref{coref2}.

\ex. \label{coref} \a. \label{coref1} What happened to my sandwich?
\b. \label{coref2} *Fido ate.

However, examples can be provided in support of the opposite. \textcite[142-144]{groefsema1995understood} makes use of sentences such as the ones in \ref{groefsema} to argue that INCs can indeed refer to specific individuals, provided sufficient linguistic context.

\ex. \label{groefsema} \a. \label{groefsema1} John brought the sandwiches and Ann ate.
\b. \label{groefsema2} John picked up the glass of beer and drank.

Nevertheless, this account does not disrupt Fillmore's firm foundation. As explained by \textcite[527]{Eu2018}, not even in sentences like \ref{groefsema} do INCs force the identification of a specific referent. What happens, instead, is that native speakers processing INCs in a flexible context of this kind can be induced to understand the missing object as if co-referring to the previously mentioned one. Thus, INCs can grammatically dissociate the mentioned referent from the one implied by the missing object, while this possibility is not active for DNCs (which are always co-referential, regardless of the context). Considering eventualities like this, it really is no wonder that \textcite[110]{Cote1996}, with respect to implicit objects in English, spoke of "murky water" in relation to the distinction between lexically-provided information and context available via world knowledge.



\paragraph{Genre-based implicit objects: a special case of definite object drop} \labpage{recipes}

I will now discuss genre-based implicit objects, a type of DNCs whose very existence goes in favor of DNCs being virtually possible with any verb, provided it appears in a discourse context that is conducive to object omission (\textit{contra} \textcite[55]{TonelliDelmonte2011}, \textit{pro} \textcite{Goldberg2001}). To quote \textcite[175]{RuppenhoferMichaelis2010}, "argument omission can but need not be licensed by a lexeme". This possibility was first observed by \textcite[95]{Fillmore1986}, who acknowledged that in "certain kinds of highly restricted mini-genres" (e.g. instructional imperatives in recipes) the omission of objects and other non-subject complements is not lexically determined (see also \textcite[237]{Haegeman1987}). Crucially, it is not the case that there is a special grammar of recipe contexts that supersedes the actual grammar of the language the recipe uses \parencite{Culy1996, Cote1996}. On the contrary, recipes and other specialized genres just serve to provide an encompassing discourse and world-knowledge context to the listener/reader. Not only that, but genres are so intertwined with argument omission that it is sometimes possible to evoke a genre just by performing the right kind of DNC, as \textcite[159]{RuppenhoferMichaelis2010} exemplify by making reference to the title of a novel by Cynthia P. Lawrence, "Chill $\varnothing$ before Serving $\varnothing$: A Mystery Novel for Food Lovers" (a clear reference to instructional imperatives found in recipes).\\
Many linguistic analyses of DNCs licensed in "mini-genres" focus on recipe contexts \parencite{Ahringberg2015, Garcia-VelascoMunoz2002, Megitt2019, Ruda2014, MassamRoberge1989, Bender, Culy1996}. In particular, \textcite{Culy1996} performed a multiple regression analysis on diachronic sets of contemporary and historical recipes with several predictors, finding that the style of a recipe book and discourse factors are the most important predictors of recipe DNCs.\\
Other authors provided accounts pertaining to a broader spectrum of genres. For instance, in addition to recipes, \textcite{Cote1996} also considers "telegraphese", i.e. the telegraphic register used in telegraphs, memos, and signs. \textcite{Weir2017} focuses on what he calls "reduced written register" in English, i.e. the absence of objects in recipes, instructional/directive imperatives\sidenote{As noted by \textcite[162]{RuppenhoferMichaelis2010}, non-instructional imperatives cannot participate in DNC constructions, as shown in their example \textit{Take *(the money) and run}.}, diaries, text messages, internet-based communication, and similar contexts. The presence of DNCs in text messages and internet-based communication is further explored by \textcite{StarkMeier2017} and \textcite{Stuntebeck2018}, focusing on Whatsapp messages. \textcite[304]{Liu2008} mentions instructional languages, such as that found on manuals, warning signs, and product labels. \textcite{Paesani2006} provides a thorough account of object (and subject) drop in special registers (such as recipes, diaries, and headlines), noting clear similarities between DNCs in recipes and the broader phenomenon of Topic drop \parencite[165]{Paesani2006}. Focusing instead on non-contemporary language found in historical texts, it is possible to find studies by \textcite{Almeida2009} on Middle English medical texts, and by \textcite{Korkiakangas2018} on Early Medieval documentary Latin.\\
In an unconventional account of football language, \textcite{BerghOhlander2016} argue that verbs licensing DNCs are "monotransitives" \parencite[54]{quirk1985grammar} in this sublanguage, since they can only take one argument. Let us consider the examples in \ref{berghohlander}.

\ex. \label{berghohlander} \a. \label{berghohlander1} Iniesta passed (the ball) and Messi finished (the attack) clinically.
\b. \label{berghohlander2} John passed *(the salt) and finished *(his steak).

In \ref{berghohlander1}, the direct objects can be omitted because the two footballers are performing acts that need no further explanation in the football community. In this game, you can only pass balls and finish attacks. Moreover, given the presence of a single ball against many players, match reports are much more likely to DNC the ball rather than the footballers, as observed by \textcite[167]{RuppenhoferMichaelis2010} and \textcite{ebeling2021score}. On the contrary, in the probable context of a dinner in \ref{berghohlander2}, it is not possible to say that John just "passed" or "finished", let alone clinically. \textcite[22]{BerghOhlander2016} explain the existence of DNCs in football reports (and, more generally, object omission) as a manifestation of the "principle of least effort" \parencite{zipf1949leasteffort} and also of the Gricean pragmatic maxim of quantity, which compels speakers to avoid being more informative than necessary. However, as \textcite[166]{RuppenhoferMichaelis2010} observed before, "genre-based omissions are never obligatory", since the maxim of quantity (favoring implicit objects) is counterbalanced by the need for informativeness (favoring overt objects).


\subsection{Neither definite nor indefinite: continuous accounts} \labsec{theory_continuous}

The account of genre-determined DNCs offered in \refsec{theory_discrete} opens the door to a broader discussion of \textcite{Fillmore1986}'s distinction between definite and indefinite omitted objects. As argued in \textcite[165]{RuppenhoferMichaelis2010} and \textcite[24]{BerghOhlander2016}, the main factor allowing for an object to be omitted is its recoverability (refer to \refsec{recoverability} for a full discussion), which depends on linguistic aspects as well as on contextual and discourse factors, and on world knowledge too. Focusing on recoverability makes it possible to go beyond the binary distinction between INCs and DNCs provided by \textcite{Fillmore1986} and many others, and also beyond the need to postulate verb-specific object-dropping capabilities. In particular, it paves the way for a non-binary account of object drop, where no clear-cut distinction between two types of omission have to be postulated (something that, in essence, was already thinkable under \textcite{HopperThompson1980}'s assumptions).\\
If recoverability is the cornerstone of object-droppability, and if it is a scalar, or even graded, feature of objects, then it stands to reason that object-droppability itself is a graded phenomenon. Taking a small step forward in this direction, \textcite{AnderBois} posits the existence of "flexible implicit arguments" to explain sentences like \textit{The Giants won $\varnothing$}, whose implicit object has a referent known to the reader (as with DNCs), which is, crucially, known because of world knowledge\sidenote{In this case, world knowledge about American football.} and not because of the presence of a linguistic antecedent (as with INCs). \textcite{CumminsRoberge2005} provide what they call a "modular account" of null objects in French, stemming from the intersection of several syntactic, semantic, pragmatic and discourse factors (a similar account of object drop, still abiding to Fillmore's binary distinction, is provided by \textcite{Cennamo2017}).\\
A more cogent, continuous account is offered by \textcite{Glass2013}, who acknowledges that there is "plenty of middle ground" between minimum recoverability (an object has to exist but it is unknown) and maximum recoverability (the specific identity of the object is known). In particular, she argues that, in order to be omitted, an object just has to be sufficiently recoverable for speakers to communicate felicitously, and that community- or genre-specific sublanguages are more prone to certain kinds of object omission simply because those smaller contexts favor object recoverability. Moreover, she explicitly argues against a INC-DNC distinction \parencite[1]{Glass2013}. A pioneering attempt to bring evidence in favor of the intuition that recoverability is the key in object omission is found in \textcite{Resnik1993, Resnik1996}, an information-theoretic account of selectional constraints testing, among other things, the relationship between transitivity and discourse context (more on Resnik's method in \refsec{resnik_sps}).\\
As I will argue later in \refsec{theory_workingdef}, I stay agnostic with respect to the binary, scalar, or continuous nature of object-droppability in my account of indefinite null objects. Following binary or scalar accounts, such a study would simply be a matter of considering those factors which are known to favor the emergence of INCs. On the other hand, under continuous-droppability assumptions, it would be a more complex matter of modeling both contexts and linguistic factors determining any kind of object drop, trying to position implicit null objects in a specific portion of the object-droppability spectrum.


\section{Defining the indefinite} \labsec{theory_defindefinite}

In this Section I will delve into a detailed discussion on indefinite implicit objects, the focus of this dissertation. In \refsec{theory_def_vs_indef} I reported a series of both now-classic and more recent accounts of the differences between so-called definite and indefinite null objects. Let us now comment on the nature of the latter, which were given several labels in the literature (objects of "detransitive verbs" in \textcite[46]{Yasutake1987}, "implicit objects" in \textcite{Glass2013} and \textcite[29]{PethoKardos2006}, or "pseudo-intransitive", "labile", "ambitransitive", "null complements", "understood arguments", "unspecified objects", "null instantiations" in other authors).

% \subsection{The involved parties} 

\subsection{Which verbs?} \labsec{whichverbs}

Traditionally (\textit{contra} this thesis and \textcite[55]{TonelliDelmonte2011}, among others), indefinite null objects are taken to only be possible with a restricted set of activity verbs. For instance, \textcite[510]{Rizzi1986} provides an argument in favor of indefinite object drop being "lexically governed" in English on the basis that in some pairs of semantically related verbs (e.g. \textit{to eat} and \textit{to devour}) one member of the pair allows for object drop, while the other does not. \textcite[236]{Haegeman1987} does not hesitate to define this account "convincing", and similar notes are also found in \textcite{Fillmore1986, Rice1988, Mittwoch2005, Gillon2012}. I will come back on the theory referring to specific case of these "semantic minimal pairs" in \refsec{mannerspec}. The important aspect, here, is that traditional or traditionally-leaning literature has trouble motivating indefinite implicit objects on the basis of meaning alone, but it also needs it to be lexically determined. With that said, which verbs does then the literature identify as allowing indefinite object drop?\\
First of all, the verbs under consideration drop \textit{syntactic} arguments, crucially, not \textit{semantic} ones. In other words, indefinite object drop are obligatory semantic arguments of such verbs \parencite[120]{Cote1996}, while they do not surface syntactically (more on the syntax of indefinite implicit objects in \refsec{theory_incorporation}). \textcite[134]{Jackendoff2003} specifically observes that it is quite inaccurate to say that such verbs "licence an optional argument", since this definition "conflates semantic and syntactic argument structure". He illustrates this point by comparing \textit{to eat} and \textit{to swallow} in \ref{eatswallow}. Both show identical syntactic behavior, but while it is possible to swallow without swallowing anything, it is not possible to eat without eating something.

\ex. \label{eatswallow} \a. \label{eatswallow1} Bill ate (the food).
\b. \label{eatswallow2} Bill swallowed (the food).

Another key point in traditional literature on indefinite implicit objects concerns the difference between change-of-state verbs (such as \textit{to break, to harden, to open}) and pseudo-transitive verbs (such as \textit{to eat, to write, to sweep}). Verbs belonging to the former class are prototypically transitive \parencite{HopperThompson1980, Kardos2010}, since they feature two maximally distinct arguments (a volitional Agent subject and a non-volitional Patient object), while verbs belonging to the latter class exhibit both transitive and intransitive features. Thus, only pseudo-transitives can licence indefinite null objects in this dichotomy. They also appear to be a semantically rich class, comprising verbs of creation (e.g. \textit{to cook, to write, to knit}), verbs of ingestion or consumption (e.g. \textit{to eat, to drink}), and verbs of surface contact (e.g. \textit{to sweep}). While syntactically they have their ambivalent behavior in common, semantically they share the fact that their objects all are "incremental themes"\sidenote{Please refer to \textcite[279]{RappaportHovavLevin2005} and \textcite[4]{Kardos2010} for an extensive account of incremental themes.}, a term originally proposed by \textcite{dowty1991thematic} to refer to verbs showing homomorphism between the physical extent of their object and the temporal progress of the event. Thus, \textit{to eat} is an incremental-theme verb because the Patient gets progressively smaller while ingested by the Agent, \textit{to write} because the Theme gets progressively more wordy while the Agent types or pens it, and so on. On the contrary, \textit{pace} Dowty's attempt to apply this analysis to change-of-state verbs \parencite[568]{dowty1991thematic}, \textit{to close} is not an incremental-theme verb because sentences like \textit{Matt closed the door half-way} do not entail that half the door was closed \parencite[279]{RappaportHovavLevin2005}. Since incremental-theme verbs can behave both transitively and intransitively, in \textcite[33]{Levin1993} they are said to participate in the "unspecified object alternation". The author also provides a list of more than 40 verbs allowing for indefinite object drop, an event which \textcite[116]{Dvorak2017} praises as a major breakthrough after previous literature only focusing, "somewhat disturbingly", on the sole verb \textit{to eat}\sidenote{An unfortunate choice for a single example of the whole category of object-dropping verbs, as we will see later in \refsec{theory_incorporation}.}.\\
As I will show with the probabilistic model of indefinite null objects I define in this dissertation (final results presented in \refch{model}), object drop is possible both with change-of-state verbs and with incremental-theme verbs, the difference being a matter of degrees (determined by several linguistic factors, see \refch{factors} and \refch{predictors}), not a binary feature as in traditional accounts.


\subsection{Which objects?} \labsec{whichobjects}

While syntactically unexpressed and semantically unspecified (at least with respect to a specific entity), indefinite implicit objects of the verbs allowing them still have to refer to \textit{something}. To what, though? Is it possible to generalize the answer?

\paragraph{Omitting \textit{something}}
This \textit{something} that object-dropping verbs refer to has been interpreted quite literally in traditional literature on the issue. \textcite{KatzPostal1967integrated} and \textcite{FraserRoss1970idioms} distinguish between the deletion of \textit{it} (expressed by the constructions which \textcite{fillmore1969types, Fillmore1986} made known as "definite null complements") and the deletion of \textit{something} (expressed by indefinite null objects). A similar consideration is also found, in passing, in \textcite[30]{PethoKardos2006} and in \textcite[6]{Ahringberg2015}.\\
However, several objections can be made to this idea. Historically, the first came from \textcite{Mittwoch1982, Mittwoch2005}, who argued that the omitted object cannot be \textit{something} because "this would be incompatible with the atelic nature of the resulting sentence". I will come back to telicity, and the somewhat different approach I will embrace in the next Chapters, in \refsec{telicity}. Other authors \parencite{Marti2015, fodor1980functional, Melchin2019, Gillon2012, Dvorak2017, Lasersohn1993} argue instead that indefinite implicit objects have to be interpreted as "weak indefinites" \parencite[55]{Melchin2019} as bare masses and plurals, instead of the indefinite pronoun \textit{something}, because only the former have obligatory narrow scope with respect to other quantifiers in the sentence (a behavior shown by indefinite null objects)\sidenote{Please refer to \textcite{carnie2021syntax}, a thorough handbook of generative syntax, for a gentle introduction to scope-taking constituents and the use of logic operators in syntax.}. Let us consider the examples in \ref{scope}, taken from \parencite[55]{Melchin2019}. In \ref{scope1} only the universal quantifier in the subject can take wide scope and the implicit object has to take the lowest scope, while in \ref{scope2} either the universal quantifier in the subject or the existential quantifier in the object can take wide scope.

\ex. \label{scope} \a. \label{scope1} Everyone ate. \hfill $\forall$ > $\exists$ / \# $\exists$ > $\forall$
\b. \label{scope2} Everyone ate something. \hfill $\forall$ > $\exists$ / $\exists$ > $\forall$

Let us unpack this notation. This means that in \ref{scope1}, for every (universal quantifier $\forall$) person, there exists some edible entity that was eaten. In \ref{scope2} this interpretation is possible too, as well as the wide-scope-existential interpretation. In this second reading, for some edible item (existential quantifier $\exists$), every person ate it.\\
The observation that indefinite implicit objects can only take low scope also holds with respect to other logic operators than quantifiers. For instance, \textcite[316]{Gillon2012} provides example \ref{gillon_neg} about negation, where the operator is shown to take obligatory wide scope over the sentence.

\ex. \label{gillon_neg} Bill did not read. \hfill $\neg \exists$x Rbx / \# $\exists$x $\neg$Rbx

This means that the only possible interpretation of this sentence is that Bill did not read anything (the first proposed truth condition of the sentence), not that there exists something that Bill did not read (the second truth condition, marked as improper).\\
Therefore, indefinite implicit objects have to be interpreted as weak indefinites (bare masses or plurals), not as the pronoun \textit{something}. However, it is clear that not \textit{any} weak indefinite can be interpreted as the omitted object of a given verb. Which ones are the right ones? I am now going to provide some answers based on the literature.

\paragraph{Prototypical objects}

\textcite[122]{vanvalinlapolla1997syntax} coined the term "inherent arguments" to describe the indefinite implicit objects occurring with activity verbs, based on the idea that they denote a facet of the meaning of the verb, characterizing the action itself rather than a participant. I will devote some space to a detailed discussion of intransitivization as a means to focus on the activity itself on \refpage{activityfocus}. Resorting once again to the concept of linguistic prototype, which is indeed central in this discussion of transitivity (\refsec{theory_transitivity}) and object drop, much literature agrees on indefinite null objects being understood as prototypical\sidenote{Note that while all referenced works appeal to the notion of prototypical argument, not all of them phrase this idea exactly in these terms. Some refer to "implied arguments", "default interpretations", "standard objects", "understood object", and other phrases in the same vein.} arguments of the verb \parencite{Rice1988, Naess2007, bresnan1978realistic, Melchin2019, Mittwoch2005, Dvorak2017thesis, Levin1993, Lorenzetti2008, quirk1985grammar}. However, what is a prototypical object of a given transitive verb? \textcite[204]{Rice1988} provides the examples in \ref{rice}.

\ex. \label{rice} \a. \label{rice1} John smokes (cigarettes / *Marlboros / *a pipe / *SMOKING MATERIALS).
\b. \label{rice2} John drinks (alcohol / *gin / *water / *coffee / *LIQUIDS).
\c. \label{rice3} When he goes to Boston, John drives (a car / *a Toyota / *a motorcycle / *A VEHICLE).
\d. \label{rice4} Each afternoon, John reads (a book / *Ulysses / *the newspaper / *PRINTED MATTER).

Examples in \ref{rice1}, \ref{rice3}, and \ref{rice4} all appeal to our world knowledge, in particular, to our knowledge of what is the most probable choice of the average Joe (or John, in these examples). People usually smoke cigarettes which do not have to be necessarily Marlboros, they drive differently branded cars in their trips out of town, and they like to read generic books in the afternoon\sidenote{Apparently, for many people, especially the ones still liking their news to be printed on paper, reading the news is a leisurely activity to be specifically enjoyed in the mornings while having breakfast. This habit had to be even more common in the late '80s than today.}.\\
As noted by \textcite[125]{Naess2007}, however, example \ref{rice2} poses a challenge to the protypicality-enabled omission theory. The most typical liquid one can drink is usually water, not alcohol. And if one was to understand that the omitted substance is alcohol due to its very omission, would not this argumentation become circular?\sidenote{See also \textcite[20]{Mittwoch2005} on the issue of circularity.} The problem alcohol poses for linguists (or better, for linguistic theory) can be explained from different angles. I, for one, would appeal to the gricean maxim of relevance, in that water is indeed the typical liquid we drink, but it is so much typical (being necessary for good health and even life) that it would be actually weird to mention water-drinking in casual conversation. No one would bat an eye at John drinking water, so it would make little sense to utter \ref{rice2} implying water-drinking. Indeed, we only refer to the act of drinking water when it becomes relevant, for instance during hot summers (\textit{Remember to drink!}) or on a Sunday morning (\textit{I am so glad I drank water before going to bed.}). Thus, since the only socially relevant, statistically likely, choice of a drink for John in \ref{rice2} is alcohol, that is what we intend as a prototypical, omissible object for the verb \textit{to drink}. This perspective is also echoed by \textcite[14]{NewmanRice2006}, who ascribe the intransitive use of \textit{to drink} to "the prominence of alcohol consumption as a topic of discourse". Another possible account is the one by \textcite[21-28]{Goldberg2005}, where "taboo verbs" are argued to facilitate object drop due to our culturally-induced shame in mentioning that which is perceived as unmentionable in polite society (such as bodily fluids or, in this case, enjoying alcoholic drinks). \labpage{drinkingdrop} Changing perspective, when presented with puzzling verb behavior such as the one expressed in \ref{rice2}, \textcite[303-305]{HuddlestonEtAl2002}\sidenote{Following \textcite[96-97]{Fillmore1986}, as discussed in \refsec{theory_def_vs_indef}.} needlessly assume that such verbs participate in two different patterns of object-droppability, i.e. "specific category indefinites" (where the omitted liquid would be interpreted as being of the alcoholic variety) and "normal category indefinites" (where the omitted liquid would be interpreted as being water). This account is flimsy at best, since it puts labels on a given state of affairs without actually providing an explanation for this epiphenomenal dichotomy. \textcite[141]{Naess2007} offers a much more compelling account, which is both descriptive and explanatory, where the "specialized meanings" \textit{to drink} and other ingestion verbs take get explained by the affectedness of their Agent\sidenote{More on Agent affectedness on \refpage{affectedagent}.}. In other words, omitting the object highlights the effect the action described by the verb has on the Agent (e.g. getting the Agent inebriated) by backgrounding the effect it has on the Patient. The fact that intransitive \textit{to drink} elicits a drink-alcohol reading much more readily than intransitive \textit{to eat} elicits a eat-a-meal reading (also noted by \textcite[14]{NewmanRice2006}) is explained by the author (and again later in \textcite[420]{Naess2011}) with reference to world knowledge and social norms. In particular, the case is made that intoxication by means of alcohol not only has a direct effect on the imbibing Agent (who takes on the drinking endeavor with this precise goal), but it also has an indirect, sometimes unintended effect, i.e. making the Agent appear visibly drunk, and thus disrespectful of several unwritten societal rules. This doubly-affected-Agent reading gives then rise to the highly specialized reading of intransitive \textit{to drink}.\\
Interestingly, \textcite[48-50]{Yasutake1987} suggests a three-way graded account of the different types of objects which can participate in indefinite implicit object constructions where the prototypicality of the omitted object is taken to be a rather flexible requirement for omission. In fact, the omissibility-as-prototypicality accounts I presented so far in this Section all made reference to the omitted object being somewhat "typical" of the verb, so that less typical Patients of the same verb are less likely to be omitted (or require quite the flight of fancy to be accounted for, as seen in the proposal by \textcite{HuddlestonEtAl2002} about the verb \textit{to drink}). Yasutake's perspective integrates the prototypicality intuition with other accounts based on world knowledge, envisioning these three types of implicit object:
\begin{itemize}
    \item typical objects (e.g. \textit{to read, to telephone});
    \item socially-understood objects (e.g. \textit{to drink, to shave, to drive});
    \item semantically unspecified objects of highly specialized activities (e.g. \textit{to steal, to see, to annihilate}).
\end{itemize}


\section{How many lexical entries?} \labsec{theory_entries}

In \refsec{theory_defindefinite} I discussed some traditional views on the characteristics a verb has to express in order to licence indefinite implicit objects and I also presented different views on the semantics of the omitted object. Now, another question arises about the nature of object-dropping verbs. When a verb participates in the so-called "implicit object alternation", is the transitive form of the verb actually distinct from the intransitive form in the lexicon, or are they different syntactic expressions of a single lexical entry? I will devote this Section to possible answers to this dilemma, to use \textcite{Gillon2012}'s word.


\subsection{Two meanings, two verbs: the naive account} \labsec{theory_twoentries}

The problem of having a single verb exhibiting two syntactically different behaviors (transitivity and intransitivity) was first identified by \textcite{fodor1980functional} and \textcite{dowty1981quantification}, an explicit reply to the former. Both treat verbs allowing for implicit objects as ambiguous between two different lexical entries, one transitive and one intransitive. This view is shared by other traditional literature on the matter \parencite{Cote1996, Mittwoch1982, vanvalinlapolla1997syntax, brisson1994licensing, FellbaumKegl1989taxonomic} and by more recent accounts \parencite{PethoKardos2006, BourmayanRecanati2013}.\\
Such an interpretation is clean on the surface, as clear-cut binarisms often are, but it does little to describe the complexity of reality \textemdash again, as binarisms often do. In a broader theory of semantics, the problem of the two uses of a single verb mirrors the well-known problem of deciding, for instance, whether \textit{bank} is a polysemous noun with two interpretations ("financial institution" and "river bank") or whether it has two homonymic interpretations. Say we go for the second, safer, account, since the only factor keeping the two senses together, i.e. etymology, is not transparent to native speakers of English nowadays. On the other hand, we would be much more keen to ascribe a polysemous interpretation to the different senses of the noun \textit{man}, which depending on the context can be used to mean "human being", "male human being", or "adult male human being". Crucially, the different senses of \textit{man} are all facets of the same entity, while the different meanings of \textit{bank} are not. Going back to the issue at hand, i.e. the distinction between transitive and intransitive senses of a given verb, it would indeed seem that these senses capture different facets of the same action performed by the Agent, instead of being two totally different meanings. This interpretation, fully consistent with the hypothesis that transitivity is a prototype (refer back to \refsec{theory_transitivity}), is further explored in \refsec{theory_incorporation} with reference to relevant literature.


\subsection{One verb, two meanings: the state-of-the-art account} \labsec{theory_incorporation}

As just argued at the end of \refsec{theory_twoentries}, transitive verbs admitting object drop are better interpreted as a single lexical entry with two different meanings, rather than two separate entries in the lexicon (an old-fashioned perspective that \textcite[60]{Lorenzetti2008} defines "counterintuitive and inappropriate"). In particular, far from being a "lexical quirk" of a restricted class of verbs, indefinite object drop appears to be "a syntactic detransitivisation mechanism" used to express events which do not embody the transitive prototype \parencite[134]{Naess2007}. Let us now discuss this behavior in more detail.

\paragraph{The syntax of indefinite null objects} First of all, a brief syntactic detour is in order. While this dissertation is much more concerned with the effect that semantics (and, to a lesser extent, pragmatics) has on indefinite object drop, it is still important to take a position with respect to the syntactic nature of the omitted object. Is it absent from the syntax, as many used to argue (as seen at the beginning of this Section and later on \refpage{activityfocus})? Or is there a syntactic slot available for the omitted object, even though it has no phonological representation? Convincing arguments brought forth by literature on the matter, as shown throughout this Section, make a strong case for the second hypothesis.\\
In syntactic theory, \textcite{roberge2002transitivity} (further explored by \textcite{CumminsRoberge2004, CumminsRoberge2005}) proposed an internal-argument equivalent of what the Extended Projection Principle (EPP) \parencite{chomsky1982epp} is for subjects, called "Transitivity Requirement". The EPP is the requirement for a subject position in the clause\sidenote{The subject position is spec-TP in languages where the subject raises to spec-TP after getting base-generated in spec-VP, or spec-VP in languages where no raising happens.}, provided by Universal Grammar, which then gets filled in by lexical material in non-\textit{pro}-drop languages (such as English) or by the empty category \textit{pro} in \textit{pro}-drop languages (such as Italian). Likewise, the Transitivity Requirement posits a direct object position in the clause, provided by Universal Grammar too, which accounts for implicit objects just like EPP accounts for null subjects. The only difference between the two is that, as noted by \textcite{CumminsRoberge2004}, "recoverability for the EPP is morphologically based, as is evident in null-subject languages, while recoverability involving the TR may also be semantically and pragmatically based". Null-subject languages have "morphological recoverability" for their null subjects in that the required information is stored in verb morphology, as shown in my examples in \ref{epp_intro} relative to Italian.

\ex. \label{epp_intro} \a. \label{epp_intro1} Corr-o. \\ run.1.SG.PRS \\ 'I run.'
\b. \label{epp_intro2} Corr-iamo. \\ run.1.PL.PRS \\ 'We run.'

Semantic and pragmatic recoverability of the implicit object licensed by the Transitivity Requirement has been partially shown in \refsec{whichobjects} and will be further discussed in the rest of this Section.


\paragraph{Focus on the activity} \labpage{activityfocus}

Many authors \parencite{Liu2008, Garcia-VelascoMunoz2002, Fillmore1986, Ahringberg2015, Levin1993, Yasutake1987, Goldberg2005} argue that the indefinite implicit object construction is used to focus on the activity itself, "downgrading the referential status of the object" \parencite[7-8]{Garcia-VelascoMunoz2002}. I provided an example of this on \refpage{drinkingdrop}, where the verb \textit{to drink} used intransitively was shown to refer to the habit of drinking alcoholic beverages. Thus, the focus of such utterances is not on the actual drink the subject is imbibing, but rather on the activity itself.\\
The word "focus" is not used idly here. Indeed the distinction between topic (the known, background information) and focus (the new, foreground information), central in pragmatic and discourse-oriented accounts of human language, also applies to the problem at hand. As argued by \textcite[66]{Lorenzetti2008}, given that every sentence requires at least one focus, and that the focus is by its very nature new, pragmatically non-recoverable information, it stands to reason that omitted objects (which are recoverable\sidenote{Refer to \refsec{recoverability}.} and prototypical\sidenote{Refer to \refsec{whichobjects}.}, hence, known) cannot be the focus. Thus, the focus in such utterances has to be on the activity itself, as in her examples in \ref{lorenzetti} (where the omitted object is \textit{people}).

\ex. \label{lorenzetti} \a. \label{lorenzetti1} I thought you said your dog doesn’t bite $\varnothing$!
\b. \label{lorenzetti2} Religion integrates $\varnothing$ and unifies $\varnothing$.

\textcite{Goldberg2005} formalized this intuition via her Principle of Omission under Low Discourse Prominence, which states that the Patient argument of a transitive verb is possible when it is "de-emphasized/unprofiled in the discourse" (i.e. neither topical nor focal) and when the action, on the other hand, is "particularly emphasised". This shift in meaning, granted by the omission of the direct object, has been shown \parencite[507]{GreeneResnik2009more} to trigger a reduced sentiment response in native speakers presented with pairs of sentences like the one in \ref{greene}.

\ex. \label{greene} \a. \label{greene1} At the same time, we should never ignore the risks of allowing the inmate to kill again.
\b. \label{greene2} At the same time, we should never ignore the risks of allowing the inmate to kill someone again.

Some authors \parencite{groefsema1995understood, recanati2002unarticulated, WilsonSperber2000, iten2005null, hall2009free} even took advantage of this focus-on-the-activity interpretation to state that the omitted object is absent in the syntax just as it is absent phonologically, being instead pragmatically provided. \labpage{activityfocus_end} Such an account is tempting and not inconsistent with Goldberg's omission principle, but (as I am going to argue right away) it is not convincing in the light of further linguistic evidence. 


\paragraph{Indefinite object drop as noun incorporation}

Not only do (indefinite) object-dropping transitive verbs describe activities, they specifically describe "conventional, name-worthy, institutionalized, habitual activities" \parencite[119]{Dvorak2017thesis}. As observed by \textcite{Marti2010, Marti2015, Yasutake1987, BourmayanRecanati2013}, this is exactly the case of verbs having undergone noun incorporation\sidenote{Noun incorporation is found in several polysynthetic languages, but it is not a requisite for polysynthesis, just as polysynthesis is not a requisite for noun incorporation \parencite{mithun2009polysynthesis}.}, a linguistic process "traditionally understood as the compounding of a noun stem with a verb stem to form a new verb stem" \parencite[5]{mithun2009polysynthesis}. Syntactically, this can be seen as a kind of head-to-head movement, as depicted in \reffig{nounincorporation_tree} (taken from \textcite[495]{carnie2021syntax}, a simplified account of the account provided by \textcite{baker1988theta}).

\begin{figure}[htb]
\caption{Portion of syntax tree illustrating the head-to-head movement involved in noun incorporation, from \textcite[495]{carnie2021syntax}.}
\labfig{nounincorporation_tree}
    \begin{tikzpicture}
\Tree[.VP [.V' [. \node(subj3){V}; ] [.DP \edge[roof]; \node(subj2){N}; ] ] ]
\draw[semithick, <-] (subj3) to [bend right=70] (subj2);
    \end{tikzpicture}
\end{figure}

Without bringing systematically noun-incorporating languages into the discussion, such as the ones from the Iroquoian family \parencite{mithun2009polysynthesis}, Frisian, or West Greenlandic \parencite{Marti2015}, it is possible to find such behavior in now-lexicalized object-verb compounds in English too. I illustrate the case of \textit{to babysit} (also valid for \textit{to birdwatch, to fingerprint}, and other such compounds) in \reffig{babysit_tree}.

\begin{figure}[htb]
\caption{Portion of syntax tree illustrating the object-verb compound \textit{to babysit} in English as a result of noun incorporation.}
\labfig{babysit_tree}
    \begin{tikzpicture}
\Tree[.VP [.V' [.V [\node(subj3){N\\baby\textsubscript{i}};] [.V\\sit ] ] [.DP \edge[roof]; \node(subj2){N\\t\textsubscript{i}}; ] ] ]
\draw[semithick, <-] (subj3) to [bend right=70] (subj2);
    \end{tikzpicture}
\end{figure}

In \textcite{Marti2015}'s account, then, the only difference between indefinite null objects and incorporated nouns such as \textit{baby} in \textit{to babysit} would be that the former are phonologically null, while the latter are not. Noun incorporation gives rise to several effects, which \textcite[455-456]{Marti2015} reports based on evidence from West Greenlandic (an ergative language):
\begin{itemize}
    \item incorporated nouns must be bare, i.e. with no preceding article/determiner (just like indefinite null objects were shown to be "weak indefinites" in \refsec{whichobjects});
    \item the subject of a noun-incorporating verb is marked with absolutive case like the subjects of intransitive (unergative) verbs and the objects of transitive verbs, while the subjects of transitive verbs are marked with ergative case (mirroring the intransitive behavior of object-dropping transitive verbs);
    \item incorporated nouns always precede the verb in the linear word order, regardless of the position a full-fledged direct object would take in the sentence;
    \item incorporated nouns are interpreted indefinitely and non-specifically (just like indefinite null objects);
    \item incorporated nouns take narrow scope with respect to other operators in the sentence (as argued about indefinite null objects in \refsec{whichobjects});
    \item verbs undergoing noun incorporation usually refer to "name-worthy, typical activities" (like indefinite null objects, as shown on \refpage{activityfocus});
    \item verbs undergoing noun incorporation tend to have "conventionalized meanings", as noted before on \refpage{drinkingdrop} relatively to the imbibe-alcohol conventionalized meaning of the verb \textit{to drink} used intransitively.
\end{itemize}

\textcite[461-463]{Marti2015} then proceeds to test the hypothesis that verbs allowing indefinite object drop and verbs allowing noun incorporation share common properties with evidence from Frisian, a Germanic language having both indefinite null objects and noun incorporation. Her analysis demonstrates that transitive-made-intransitive verbs and incorporated-into verbs indeed belong to the same class. A crucial property both types of verbs share is, for instance, that only verbs selecting for a Patient object and having a volitional subject can participate in these constructions (making it only possible for verbs such as \textit{to notice, to hate, to know} in English to participate in \textit{definite}, not \textit{indefinite}, object drop).\\
Purely pragmatic accounts of indefinite implicit objects, as revealed on \refpage{activityfocus_end}, fail to take into account the cluster of properties shared by noun-incorporating and object-dropping verbs alike. 


\paragraph{Affected Agents, effected Patients} 

I will now tackle an aspect of indefinite objecthood originating directly from the account of transitivity as a prototype concept by \textcite{HopperThompson1980} and later literature\sidenote{Refer back to \reftab{ht1980_parameters} and further considerations in \refsec{theory_transitivity}.}, i.e. the need for the subject and the object of transitive sentences to be maximally distinct in their semantic behavior. This requirement was formalized by \textcite[30]{Naess2007} in the Maximally Distinguished Arguments Hypothesis. This observation led some authors\sidenote{Refer back to \refsec{whichverbs}.} to posit change-of-state verbs such as \textit{to break} as obligatorily transitive, since they feature a volitional, unaffected Agent and a non-volitional, affected Patient. As argued time and again in this Chapter, indefinite object drop is far from being prototypical behavior for transitive verbs. Under this lens, then, it would make sense to find that indefinite null objects are more common with verbs having affected Agents (i.e. Agents being the endpoint of the event) and/or effected Patients (i.e. Patients brought into existence by the event the verb refers to). Such an analysis is discussed in detail in \textcite{Naess2007}.\\
\labpage{affectedagent}Let us consider affected Agents first. \textcite[158]{tenny1994aspectual} calls them "measuring arguments", in that they delimit the event "by undergoing a change of state that marks the temporal end of the event". For instance, the event described by \textit{John ate} would be delimited by the affectedness of the Agent (in this case, the feeling of fullness), rather than by the affectedness of the unmentioned Patient. On the opposite, if one wanted to focus on the affectedness of the Patient, they would then have to resort to a transitive use of the verb \textit{to eat} \parencite[80]{Naess2007}. An affected-Agent interpretation of intransitive \textit{to eat}\sidenote{Such an account was first offered, \textit{in nuce}, in \textcite{Wierzbicka1982, starosta1978valence, nedjalkov1988typology, haspelmath1994passive}.}, which may or may not convince the reader yet by means of this example in English, becomes quite more persuasive in \textcite[61-63]{Naess2007}'s account of the same verb in Yucatec, a Mayan language spoken in the Yucatán Peninsula. In this language, due to its Agent being affected, \textit{to eat} patterns aspectually and morphologically with change-of-state verbs, not with activity verbs (as one would expect). Crosslinguistically \parencite[126]{Naess2007}, ingestive verbs (\textit{to eat}, \textit{to drink}, but also \textit{to learn}) consistently show characteristics typically belonging to intransitive verbs \parencite{Amberber2009}, leading \textcite{marantz1981nature} and subsequent decades of literature on indefinite object drop to consider them "class representatives" of the typical behavior of verbs allowing for indefinite object drop. Taking the affected-Agent interpretation to the extreme, in some languages (such as Korean and Turkish) \textit{to eat} is even used as a grammaticalized marker of Agent affectedness, e.g. as an auxiliary, as a light verb, in constructions where it expresses undergoing or adversativity \parencite[75]{Naess2007}, and in antipassive constructions\sidenote{Used in ergative languages to convey certain aspectual and modal nuances. They are called "antipassives" because, much like in passive sentences the object of the active sentence becomes the subject and the subject of the active sentence is deleted, in antipassive constructions the object is (usually) deleted and the subject changes case from ergative to absolutive case.} \parencite[414]{Naess2011}. \labpage{antipassives} \textcite{Nicolas2019} offers an interesting account of the features null objects in English share with antipassive constructions. Moreover, an affected-Agent account can be easily employed to explain linguistic behavior that would otherwise remain unmotivated, such as the resistance of the verb \textit{to lock} to intransitive use as opposed to the ease one finds in using \textit{to eat} intransitively, as noted by \textcite[30]{PethoKardos2006}. They argue that the opposite behavior of these two verbs with respect to indefinite object drop "does not become clear", on the basis of their selection restrictions\sidenote{More on selectional restrictions and object recoverability in \refsec{recoverability}.} having comparable extent (i.e. the objects of \textit{to eat} are edible items, the objects of \textit{to lock} are items provided with a lock). The reason for this difference, however, becomes quite clear when one considers that the subject of \textit{to eat} is an affected Agent, while this does not hold true for \textit{to lock}. It is also important to note that while Agent affectedness is inherent to the semantics of some verbs (such as ingestion verbs), it may also be activated by verb-external elements of a clause. Let us consider the example sentence \textit{John murdered for the money} \parencite[136]{Naess2007}. In this case, the affected-Agent interpretation is fostered by the purpose clause \textit{for the money}, since the Agent's motive for the homicide is a direct gain, i.e. something that positively affects the Agent. Finally, the affected-Agent account can also explain constructions of the type \textit{have a drink, have a lick, have a bite} \parencite[758, 771]{Wierzbicka1982}. The author argues that in such constructions, \textit{have} has a detransitivizing function in that it backgrounds the object while focusing on the Agent. Not only that, but also, \textit{have a [verb]} events deviate from the transitive prototype because their Agent is affected by it (typically, by enjoying the activity the verb refers to), while the Patient is minimally affected.\\
Let us now discuss the other kind of argument deviating from the transitive prototype, i.e. effected Patients. These are non-affected objects that come into existence thanks to the very action described by the verb, and only if this action is brought to completion, e.g. the letter in \textit{John is writing a letter} or the cake in \textit{John is baking a cake}. Such constructions tend to feature indefinite null objects crosslinguistically and to have unaffected Agents, making it necessary to provide a different analysis than before in this Section \parencite[127-128]{Naess2007}. What affected-Agent and effected-Patient constructions have in common is that they both show low distinctness between Agent and Patient (making them the optimal environment for felicitous indefinite object drop), and they both largely depend on the semantics of the verb itself. As \textcite[127]{Naess2007} and \textcite[421]{Naess2011} observe, the low-distinctness of effected Patients is so embedded in the verb semantics that the very intransitive use of an effected-Patient verb evokes the non-referentiality of the object. This is most evident in imperfective contexts\sidenote{I will come back to the role of (im)perfectivity as a factor determining indefinite object drop in \refsec{perfectivity}.} (e.g. \textit{John was writing}), where the effected Patient is not presented as fully effected yet, and thus it is even less prominent in the discourse. On the contrary, perfective contexts tend to block indefinite object drop with effected-Patient verbs (e.g. \textit{? John had written}), while the same does not usually hold for affected-Agent verbs (e.g. \textit{John had eaten}). Crucially, an effected-Patient account can be used to explain why, as first observed by \textcite[96]{Fillmore1986}, intransitive \textit{to bake} in English (e.g. \textit{I spent the afternoon baking}) can only be understood to refer to the act of baking "bread or pastries, but not potatoes or ham". \textcite[135]{Naess2007} easily explains this linguistic fact by observing that the bread-or-pastries interpretation features an effected Patient, while the potatoes-or-ham interpretation features an \textit{affected} Patient, which makes the verb prototypically transitive and, thus, resistant to indefinite object drop. On the same note, intransitive \textit{to paint} (e.g. \textit{He paints}) is interpreted to refer to the act or habit of painting pictures, not house walls.


\section{A working definition of "indefinite object drop"} \labsec{theory_workingdef}

This dissertation is about indefinite null objects, whose nature I discussed throughout this Chapter with reference to traditional and recent literature alike. It is now time to end this Chapter with a note on some more detail on what this thesis is going to focus on, and also on what is going to be ignored in my experimental account of indefinite object drop.\\
As I mentioned in \refsec{theory_continuous}, I assume there is a distinction between \textit{indefinite} object drop, motivated by linguistic factors (see \refch{factors} and \refch{predictors}) of varying nature, and other kinds of object drop, such as definite null objects depending on Topic drop, genre-based null objects, and implicit objects depending on extra-linguistic context. This choice is motivated by lexical, semantic, and pragmatic accounts of the indefinite object drop discussed in this Chapter. Crucially, while I assume that verbs participating in the indefinite object drop construction show two surface sentence structure while being single lexical entries (\refsec{theory_entries}), I will not take a stance with respect to the problem of whether indefinite and definite object drop form part of a continuum or are two discrete, binary phenomena (\refsec{theory_def_vs_indef}). Indeed, one could interpret the distinction between the different kinds of object drop as stemming from a situation like the one in \reffig{objectdrop_venndiagram}, where virtually any verb can participate in object-dropping constructions, but in different ways depending on the underlying factor.

\begin{figure}[htb]
\caption{A concentric view of the factors determining the possible continuum between definite and indefinite object drop.}
\labfig{objectdrop_venndiagram}
    \begin{tikzpicture}
\coordinate (O) at (0,0);
\foreach \j in {1,...,3} \draw (O) circle (3.5-\j);
\foreach \k/\text in {0/extra-linguistic context,1/intra-linguistic context,2/semantics} \draw[decoration={text along path,reverse path,text align={align=center},text={\text}},decorate] (2.6-\k,0) arc (0:180:2.6-\k);
    \end{tikzpicture}
\end{figure}

Let us interpret the examples in \ref{concentric} in the light of \reffig{objectdrop_venndiagram}. In \ref{concentric1}, the referent of the omitted object (gas) is supplied by pragmatics \textemdash if one wanted to echo \textcite{Glass2020}, it is recoverable because in the sub-community of car-drivers it is customary to say cars "drink" to refer to them burning fuel. In \ref{concentric2}, the definite null object is immediately recoverable from the linguistic context provided in the sentence, and no additional extra-linguistic context is needed to interpret this utterance. Instead, the indefinite null object in \ref{concentric3} is interpretable with no context whatsoever, given that the semantics of the verb (and of the Agent, as argued on \refpage{affectedagent}) are sufficient for its recoverability.

\ex. \label{concentric} \a. \label{concentric1} It drinks $\varnothing$ a lot! \\ (in the social context of someone speaking about a car)
\b. \label{concentric2} \# My milk has been opened, who drank $\varnothing$?
\c. \label{concentric3} He drinks $\varnothing$.

One could, naturally, have qualms with respect to the full grammaticality of \ref{concentric2}, since English, after all, is not typically a language allowing definite null objects outside genre-specific environments. Moreover, one may want to keep verb-specific semantic affairs separate from extra- and intra-linguistic contextual information. In this case, the concentric, continuous view depicted in \reffig{objectdrop_venndiagram} would not correspond anymore to the theory, and one would have to adopt a strictly binary definite-or-indefinite perspective. Such an account easily handles cases where the same verb can appear in indefinite object drop constructions when proper situational context is provided, but not in zero-context environments, as in \ref{bindrop}.

\ex. \label{bindrop} \a. \label{bindrop1} Do you even lift $\varnothing$, bro? \hfill \parencite[9]{Glass2020} \\ (common among people training for strength in gyms)
\b. \label{bindrop2} *John had lifted $\varnothing$.

What is important, here, is that the same factors acting on no-context indefinite null objects are also active in context-rich scenarios, regardless of the continuous or binary view one adopts. What matters is that additional context (such as the one granted by a sub-genre) makes it possible for a verb to license indefinite implicit objects more freely than in no-context utterances. Thus, by studying the factors determining indefinite object drop in no-context utterances, I am still providing useful information to understand context-depending object drops, although without committing to a specific interpretation of the distinction between indefinite implicit objects and other null objects.
